Das III. Physikalische Institut der RWTH Aachen University betreibt einen kosmischen Teststand mit einer voll funktionsfähigen \ac{DT} Kammer und der zugehörigen Peripherie. Die Driftkammer entspricht den verwendeten Kammern im \ac{CMS} Experiment. In diesem Jahr wurde die Kammer mit der Ausleseelektronik, die im Phase-2 Upgrade von \ac{CMS} verwendet wird, ausgestattet und in Betrieb genommen. Außerdem wurde am Institut ein Szintillationsdetektor mit Silizium-Photomultiplier (\acs{SiPM}) Auslese entwickelt. Mit dem Ziel der Integration der Auslese des Szintillationsdetektors in die Backend-Elektronik der Driftkammer wurde neue Szintillator-Ausleseelektronik und -Firmware entwickelt. Diese nutzt den \ac{lpGBT} chip, ein Serializer-Deserializer Modul, welches am \acs{CERN} entwickelt wurde. Der Chip erlaubt die Verbindung des Auslesesystems zur \ac{DT} Backend-Elektronik und die Synchronisation mit der globalen Timing-Referenz. Diese Arbeit präsentiert die neu entwickelte Auslesekette für den Szintillator und erste aufgenommene Daten. Die Daten werden mit Teilchenspuren der \ac{DT} Kammer verglichen, welche mit einem selbst implementierten, triggerlosen Rekonstruktionsalgorithmus erhalten werden. Der Algorithmus basiert auf der lokalen \ac{DT} Rekonstruktion, welche im \ac{CMS} Experiment verwendet wird. Die Synchronization der beiden Detektoren konnte in dieser Arbeit erreicht werden, mit einer Zeitpräzision in der Größenordnung von $\mathcal{O}(\SI{10}{\nano\second})$. Die implementierte Spur-Rekonstruktion konnte erfolgreich validiert werden, indem die Szintillator-Hits als Referenz genutzt wurden.

\acresetall